\section{Diskussion}
Die Ergebnisse dieser Arbeit verdeutlichen, dass die wahrgenommene Bequemlichkeit ein zentraler Treiber für die verstärkte Nutzung des Online-Shoppings während der Corona-Pandemie war. Dieses Ergebnis steht dabei im Einklang mit den bestehenden Annahmen der Konsumforschung, wonach die Convenience ein wesentlicher Stimulus für digitale Kaufentscheidung ist.

Im Bezug auf das S-O-R-Modell, kann hier die Pandemie als externer Stimulus angesehen werden, der die bestehenden Routinen disruptiert und neue Verhaltensweisen angestoßen hat. Die wiederholten Onlinekäufe während der Pandemie entwickelten sich bei mehreren Befragten zu stabilen Gewohnheiten.

Zusammenfassend lässt sich sagen, dass die Corona-Pandemie das individuelle Kaufverhalten der Befragten deutlich in die Richtung des Online-Shoppings verschoben hat. Diese Verschiebung entstand vor allem durch die wahrgenommene Bequemlichkeit, erhöhte Vergleichsmöglichkeiten sowie situative Einschränkungen während der Pandemie.

Gleichzeitig zeigen die Ergebnisse, dass keine vollständige Verdrängung des stationären Einzelhandels stattgefunden hat. Hier hatte sich bei Befragten ein hybrides Kaufverhalten etabliert, bei dem die Wahl zwischen Online-Shopping und dem stationären Handel zunehmend von situativen Faktoren abhängig war.

Allerdings sind aus den Ergebnissen der vorliegenden Arbeit auch Limitationen ersichtlich. Die Studie basierte auf einer kleineren Stichprobe von nur drei Personen, wodurch eine Übertragbarkeit der Ergebnisse eingeschränkt ist. Außerdem können bei den Befragten retrospektive Verzerrungen der Erinnerung auftreten, welche hier nicht ausgeschlossen werden konnten. Außerdem ist zu berücksichtigen, dass die Interviews einen eher geringen Umfang hatten und sehr kurz ausfielen, was möglicherweise nicht alle relevanten Aspekte des Kaufverhaltens ausreichend beleuchtete.

Für den stationären Einzelhandel legen die Ergebnisse nahe, dass besonders das Einkaufserlebnis und die haptische Produkterfahrung weiterhin ein sehr wichtiger Faktor im Verleich zum Online-Shopping bleiben wird.

Zukünftige Forschungen in diesem Bereich sollten sich mit größeren und vielfältigeren Stichproben beschäftigen, um zu prüfen, ob die beobachteten Gewohnheitsverschiebungen langfristig weiterbestehen.